\section{Branching Implementation Details}
\label{app:branching}

The core insight behind our training pipeline is that debate trajectories
are not linear sequences but \emph{trees}: at key decision points,
multiple viable continuations exist, and the opponent's response may differ
depending on which path is taken. We exploit this observation to generate
high-signal counterfactual training pairs.

\subsection{BranchDecider Algorithm}

At each speech point, the multi-trial generation system produces $k$
candidate speeches (typically $k{=}4$: two temperature variants and two
retry variants). The \textsc{BranchDecider} determines whether any pair
of candidates would produce a meaningfully different debate continuation.

\paragraph{Divergence score computation.}
For each pair of candidate speeches $(s_i, s_j)$, the BranchDecider
computes a composite divergence score across four dimensions:

\begin{table}[h]
\centering
\small
\begin{tabular}{lp{7cm}c}
\toprule
\textbf{Divergence Type} & \textbf{Criteria} & \textbf{Weight} \\
\midrule
Tactical & Different macro-tactic or substantially different micro-tactic sequence & 0.30 \\
Structural & Different contention order, dropped vs.\ addressed arguments, or different voting issues & 0.25 \\
Evidential & Different evidence cards cited or substantially different warrant construction & 0.25 \\
Rhetorical & Different framing, tone, or persuasive strategy (e.g., pathos vs.\ logos emphasis) & 0.20 \\
\bottomrule
\end{tabular}
\caption{Divergence dimensions and their weights in the BranchDecider scoring.}
\label{tab:divergence_types}
\end{table}

The composite divergence score is the weighted sum:
\begin{equation}
d(s_i, s_j) = \sum_{t \in \{\text{tac}, \text{str}, \text{evi}, \text{rhe}\}} w_t \cdot d_t(s_i, s_j)
\end{equation}

\noindent where each $d_t \in [0, 1]$ is assessed by the scoring LLM
alongside the standard quality evaluation.

\paragraph{Branching conditions.}
A branch is created when \emph{both} of the following hold:

\begin{enumerate}
\item \textbf{Divergence threshold}: $d(s_i, s_j) \geq 0.15$.
\item \textbf{Viability check}: Both candidates score $\geq 0.65$ on inline quality. This prevents branching on a high-quality speech versus a degenerate one, which would produce a trivial preference pair.
\end{enumerate}

The key question operationalizing divergence is: \emph{``Would the opponent
respond differently to $s_i$ versus $s_j$?''} A high divergence score
implies yes, meaning the downstream trajectory will differ substantively.

\paragraph{Branch limits.}
Each debate permits at most 4 active branches to bound computational cost.
When more than 4 branch points are identified, the system selects the
4 with highest divergence scores. In practice, most debates produce 1--2
branches, concentrated at the NC (Negative Constructive) and 1AR
(First Affirmative Rebuttal) speech points where strategic optionality
is highest.

\subsection{Branch State Tracking}

Each branch maintains a \texttt{BranchState} record:

\begin{table}[h]
\centering
\small
\begin{tabular}{lp{8cm}}
\toprule
\textbf{Field} & \textbf{Description} \\
\midrule
\texttt{branch\_id} & Unique identifier (format: \texttt{debate\_id.branch\_point.variant}) \\
\texttt{parent} & ID of parent branch (null for the root trajectory) \\
\texttt{child\_branches} & List of child branch IDs created from this branch \\
\texttt{speeches} & Ordered list of speeches from the branch point onward \\
\texttt{divergence\_point} & Speech index where branching occurred \\
\texttt{divergence\_score} & Composite $d(s_i, s_j)$ at the branch point \\
\texttt{outcome} & Final debate outcome (win/loss and score) for this branch \\
\bottomrule
\end{tabular}
\caption{BranchState record fields.}
\label{tab:branch_state}
\end{table}

Branches share all speeches before the divergence point and maintain
independent copies thereafter. The opponent model generates fresh
responses for each branch, ensuring that downstream effects of the
divergent choice are faithfully captured.

\subsection{Winner-Shift Detection}

The highest-signal training examples arise when different branches
produce different debate outcomes. The winner-shift detection algorithm
proceeds as follows:

\begin{enumerate}
\item For each branch pair $(B_a, B_b)$ sharing a parent, compare final outcomes.
\item A \emph{winner shift} occurs when $\text{outcome}(B_a) \neq \text{outcome}(B_b)$---i.e., the trained model wins one branch and loses the other.
\item For winner-shifted pairs, construct a \texttt{CounterfactualInsight}:
\end{enumerate}

\begin{table}[h]
\centering
\small
\begin{tabular}{lp{8cm}}
\toprule
\textbf{Field} & \textbf{Description} \\
\midrule
\texttt{losing\_choice} & The speech variant that led to the losing branch \\
\texttt{winning\_choice} & The speech variant that led to the winning branch \\
\texttt{impact\_explanation} & LLM-generated causal narrative explaining why the choice mattered \\
\texttt{confidence} & Estimated reliability of the causal attribution (0--1) \\
\texttt{score\_differential} & $|\text{score}(B_\text{win}) - \text{score}(B_\text{lose})|$ \\
\bottomrule
\end{tabular}
\caption{CounterfactualInsight structure for winner-shifted branch pairs.}
\label{tab:counterfactual_insight}
\end{table}

The \texttt{impact\_explanation} is generated by prompting the scoring LLM
with both complete branch trajectories and asking it to identify the causal
chain from the divergent choice to the outcome difference. These explanations
are preserved for potential use in reasoning distillation.

\subsection{Preference Pair Types and Weighting}

Branching produces four distinct types of GRPO preference pairs, each
weighted according to its estimated signal quality:

\begin{table}[h]
\centering
\small
\begin{tabular}{lccp{5.5cm}}
\toprule
\textbf{Pair Type} & \textbf{GRPO Weight} & \textbf{Typical \%} & \textbf{Construction} \\
\midrule
Trial pairs & $1.0\times$ & 55--65\% & Best vs.\ worst among $k$ trials at a single speech point (no branching) \\
Branch pairs & $2.0\times$ & 10--15\% & Winner-shifted branch pairs: the winning branch's speech is chosen, the losing branch's speech is rejected \\
Golden pairs & $1.5\times$ & 15--20\% & Opus-generated reference vs.\ best model trial; used in SFT pretraining and as high-quality anchors \\
Retry pairs & $1.0\times$ & 10--15\% & Original vs.\ retry when score gap $\geq 0.10$; higher-scoring variant is chosen regardless of source \\
\bottomrule
\end{tabular}
\caption{GRPO preference pair types with weights. Branch pairs receive $2\times$ weight because they carry outcome-level causal signal rather than speech-level quality signal alone.}
\label{tab:pair_types}
\end{table}

Branch pairs receive elevated weight because they encode a stronger training
signal: not merely ``this speech scored higher'' but ``this choice changed
who won the debate.'' The $2.0\times$ weight was selected via ablation
(Section~\ref{app:grpo_groups}); weights of $1.5\times$ and $3.0\times$
produced lower validation accuracy on held-out debates.
